% Introduction (owner: A).
\section{Introduction}
\label{sec:intro}

Coding agents routinely leave the editor to retrieve documentation, inspect a web page, or submit a
result.  Some use a built-in fetch service; others ask the operating system's curl or Python runtime
to make the request.  In the latter case the visited site does not merely see an ``AI agent.''  It
sees a connection produced by software installed on the user's machine, together with its TLS and
HTTP choices, the user's network address, and the delay while a model decides what to do next.  A
network of sites could therefore use ordinary requests to infer which agent and model are active and,
more consequentially, to link different agents back to the same machine or operator.

Prior work fingerprints browser agents from rich network and UI traces
\cite{kang2026whose,lugoloobi2026actions,wang2026fpagent,fayolle2026browser}, identifies models from
their generated text~\cite{pasquini2025llmmap,sun2025idiosyncrasies}, or links conventional browsing
histories across devices~\cite{zimmeck2017crossdevice,bird2020replication}.  Coding agents occupy a
different point in this space: they often execute no JavaScript, make only a few HTTP requests, and
may delegate networking to generic command-line tools.  This creates a gap between \emph{agent
identification} and \emph{user linking}.  A built-in client can identify the agent while hiding the
local environment; a shell command can be indistinguishable from a script yet expose that
environment.  We ask three questions: (1) can a site distinguish an agent from a plain program using
the same HTTP stack, (2) can it infer the model from submitted text or request behaviour, and (3) can
cooperating sites link visits across agents, sites, network changes, and the user's browser?

We answer these questions with a logging HTTPS server that records TLS, HTTP/1.1, HTTP/2, timing, and
submitted-text features.  We compare Claude Code and Codex CLI against 12 plain HTTP stacks on one
Linux and two Windows machines, then run a controlled campaign on two public hostnames.  Every run
has ground-truth labels, but those labels are excluded from the observable feature sets.  The design
uses no browser JavaScript, hidden instruction, or prompt injection: each page visibly asks the agent
to summarize and submit its contents, and all participants are the authors operating fictitious
personas on their own machines.

Our contributions are:
\begin{itemize}
  \item a reproducible, multi-layer measurement harness and same-machine script baselines for coding
  agents, separating what belongs to the agent from what belongs to the user's networking stack;
  \item evidence that shell-using agents can match ordinary scripts at the TLS and HTTP layers while
  model-in-the-loop timing and within-run tool choice still expose agent and model information; and
  \item a cross-site linking protocol that measures the stability and collisions of network and
  environment evidence across agents, browsers, machines, and changed IP addresses, while reporting
  separately the cases that cannot support a reliable link.
\end{itemize}
